\documentclass[11pt]{article}
\usepackage[margin=1in]{geometry}
\usepackage{amsmath,amssymb,amsthm,mathrsfs}
\usepackage[hidelinks]{hyperref}
\newtheorem{lemma}{Lemma}
\newtheorem{theorem}[lemma]{Theorem}
\newtheorem{corollary}[lemma]{Corollary}
\title{Scaled single-digit terminal families for ordinary gap six}
\author{Internal proof; no priority or external-review claim}
\date{18 November 2026}
\begin{document}\setlength{\emergencystretch}{1em}\maketitle
We prove ordinary Strong Factorial for $X(A+BX^6)$ at every terminal
$N=b7^r+1$, $3\le b\le6$, $r\ge0$.
Together with the two smaller leading digits this covers
$1\le b\le6$.  The argument retains every boundary row.
A uniform formal generating-series calculation supplies their divided
reductions; fixed polynomial identities and periodic finite-field
orbits then exclude rank loss at every digit depth $r$.
This is a terminal family, not all-index gap-six Strong Factorial.

For now fix $3\le b\le6$, $r\ge1$, $L=7^r$, $N=bL+1$ and $c=N+5$.
Work in a complete algebraically closed extension of $\mathbb Q_7$,
with valuation ring $\mathcal O$ and residue field $k$; normalize
$v_7(7)=1$.  Put $\mathcal A=c-zD_z+z$ and
$\Pi_N(X)=\prod_{\ell=0}^{N-1}(X+6\ell)$.
For $Q\in\mathcal O[X]$, $\deg Q\le4$, let
\[
 G_Q(X)=\binom{X/6+N-1}{N}Q(X),\qquad
 d_j(Q)=\sum_{\ell=0}^j(-1)^\ell\binom j\ell G_Q(c-\ell).
\]
Every $d_j$ is integral: the binomial polynomial is integral on
$\mathbb Z_7$ by continuity from nonnegative integers, and the
coefficients of $Q$ are integral.
Also $d_j=0$ for $j>N+4$.
The finite formal conjugacy of $\mathcal A$ gives
\[
 \frac{\Pi_N(\mathcal A)Q(\mathcal A)1}{6^NN!}
   =\sum_{j=0}^{N+4}\frac{d_j(Q)}{j!}z^j .
\]
No exponential is evaluated here.  At $u=z^6=7v$ define the six
integral polynomial rows
\[
 V_a(Q;v)=\sum_{\ell\ge0}
   \frac{7^\ell a!}{(a+6\ell)!}d_{a+6\ell}(Q)v^\ell,
                \qquad0\le a<6.                         \tag{1}
\]
Their simultaneous zero is the actual cyclic relation.
Integrality follows from the exact weight valuation
$(s_7(a+6\ell)-a)/6\ge0$, with $s_7$ the base-seven digit sum.

\begin{lemma}[The leading rows]\label{lem:digit-rows}
Put
\[
 S=S_r(v)=1+\sum_{i=1}^r(-1)^iv^{(7^i-1)/6},\quad
 f(x)=(x+1)\bar Q(6-x),\quad
 D_a=\sum_{\ell=0}^a(-1)^\ell\binom a\ell f(\ell).
\]
Then $\bar V_a=0$ for $a<b$.  For $b\le a\le5$,
\[
 \bar V_a(Q)=\binom ab\sum_{j=0}^b
                   \binom bj(-S)^jD_{a-j}.              \tag{2}
\]
In particular each entire row $V_a$, $a<b$, is coefficientwise
divisible by seven.
\end{lemma}
\begin{proof}
Only indices with $s_7(j)=a$ contribute to $\bar V_a$.
For Lucas-allowed $\ell$ digitwise bounded by such $j$, its digit
sum is at most five.
The top binomial argument in $G_Q(c-\ell)$ is
$1-\ell/6+b7^{r+1}/6$.
Since $-1/6=1+7+7^2+\cdots$, its units digit is $1+\ell_0$
and its $r$th digit is $s_7(\ell)$, without carries.
Lucas' formula therefore gives
\[
 \overline{G_Q(c-\ell)}
      =f(\ell_0)\binom{s_7(\ell)}b .
\]
Every relevant $j$ has positions at most $r$, since
$j\le N+4=b7^r+5<7^{r+1}$.
If $J=\sum_{i\ge1}j_i$, summing the higher digit differences gives
the exact expression
\[
 \bar d_j=[\xi^b](-\xi)^J
   \sum_{\ell=0}^{j_0}(-1)^\ell
                 \binom{j_0}{\ell}f(\ell)(1+\xi)^\ell.   \tag{3}
\]
It vanishes when $a<b$.
For completeness, the factorial unit is
$7^{-v_7(j!)}j!\equiv(-1)^{v_7(j!)}\prod_i j_i!$.
Thus the higher-digit weights at total higher degree $J$ sum to
$(S-1)^J/J!$.
Inserting (3), and using
$\binom{j_0}{\ell}\binom{\ell}{b-J}
 =\binom{j_0}{b-J}\binom{j_0-b+J}{\ell-b+J}$,
gives
\[
 \binom ab(-1)^b
   \sum_{J=0}^b\binom bJ(S-1)^J
             E^{b-J}(1-E)^{a-b}f(0)
 =\binom ab(1-E-S)^b(1-E)^{a-b}f(0),
\]
where $Ef(x)=f(x+1)$.
Expansion is exactly (2).
For $a=0$ the constant term also vanishes modulo seven, by the
same Lucas calculation; every other coefficient has positive
weight valuation.  This proves the whole-row divisibility.
\end{proof}

\begin{lemma}[Uniform divided-row seeds]\label{lem:divided-seeds}
For $a<b$ write $R_a=\overline{V_a/7}$.
At $Q=1$ these rows have the following values in $k[S,v]$:
\[
\begin{array}{c|l}
 b&\text{divided-row values at }Q=1\\\hline
3&(R_0,R_1,R_2)=(6,\ 4-S,\ 2+5S-S^2)\\
4&(R_0,R_1,R_2,R_3)=(1,\ 4+S,\ 6+6S+S^2,\ S+3S^2+S^3)\\
5&R_0=6+vS^5,\quad R_1=1-S,\\
 &R_2=2+6S+6S^2,\quad R_3=3S+5S^2+6S^3,\\
 &R_4=5S^2+3S^3+6S^4\\
6&R_0=1-vS^5(S-1),\quad R_1=S-2+vS^6,\\
 &R_a=S^{a-2}(1+5S+S^2),\quad 2\le a\le5
\end{array}                                                     \tag{4}
\]
\end{lemma}
\begin{proof}
All following series are formal $t$-adic, not boundary values of
analytic exponentials.
Lemma \texttt{lem:diagonal} of
\href{GAP_SIX_TWO_POWER_PLUS_ONE_SFC.tex}{the two-power-plus-one note}
gives the exact series
\[
 \sum_{m\ge0}\frac{B_{m,m,6}(z)}{6^mm!}t^m
 =\kappa e^{zH},\quad
 R-R^7=t,\ R(0)=1,\ H=1-R,\quad
 \kappa=\frac{6R}{7R^6-1}\equiv R-7t\pmod{49}.           \tag{5}
\]
The coefficient of $t^N$ in $a!\kappa\mathcal C_{6,a}e^{zH}$
is precisely $V_a(1;v)$, where $\mathcal C_{6,a}$ extracts the
coefficient of $z^a$ after $z^6=7v$.

The Artin--Hasse series
\[
 \mathcal E_7(X)=\exp\left(\sum_{i\ge0}X^{7^i}/7^i\right)
\]
has integral coefficients by
\href{ARITHMETIC_FRONTIER_TERMINAL_ORDERS.tex}
{Lemma \texttt{lem:artin-hasse-closure} of the arithmetic note}.
There is a unique integral series
$\mathscr A\in t\mathbb Z_7[v][[t]]$ satisfying
\[
 \mathscr A+v\mathscr A^7+7^6v^8\mathscr A^{49}+\cdots=H
\]
has $e^{zH}=\mathcal E_7(z\mathscr A)$ in the formal cyclic quotient.
Every fixed $t$ coefficient uses only finitely many $z$ powers.
For $1\le a\le5$ this gives
\[
 a!\mathcal C_{6,a}e^{zH}
 \equiv\mathscr A^a+7c_av\mathscr A^{a+6}\pmod{49},
 \quad c_a=a![X^{a+6}]\mathcal E_7(X)\in\mathbb Z_7.
\]
For $a=0$ it gives $1+7v\mathscr A^6/6!$.
Every subsequent section term is divisible by49, since
Artin--Hasse coefficients are integral.

Let $T=\sum_{i\ge0}t^{7^i}$, $a_0=-1$,
$a_i=-1-va_{i-1}^7$, and $A_0=\sum_{i\ge0}a_it^{7^i}$.
For $U=T$ or $A_0$ put
$C(U)=(U^7-\operatorname{Fr}(U))/7$, with
$\operatorname{Fr}(A_0)=\sum_i a_i^7t^{7^{i+1}}$.
These carry series are integral and supported at exponents1 modulo6.
The implicit equations give
\[
 R=1+T+7R_1,\quad R_1\equiv C(T)+P(T),\quad
 \mathscr A=A_0+7A_1,\quad A_1\equiv-R_1-vC(A_0)\pmod7,
\]
where $P(T)=T+3T^2+5T^3+5T^4+3T^5+T^6$.
Also $a_r\equiv-S_r(v)$.
For $a<b$, the coefficient of $t^N$ in $(1+T)A_0^a$ is
exactly zero: it uses at most $a+1\le b$ powers of seven, while
$N$ has digit sum $b+1$.
Consequently, for $1\le a<b$, division by seven gives
\[
 R_a(1)=[t^N]\{R_1A_0^{a-1}(A_0-a(1+T))
 -tA_0^a-av(1+T)A_0^{a-1}C(A_0)
 +c_av(1+T)A_0^{a+6}\}.                                \tag{6}
\]
For $a=0$ the expression is
\[
 R_0(1)=[t^N]\{C(T)+P(T)-t+6v(1+T)A_0^6\}.             \tag{7}
\]

If $b\le5$, all carry and Artin--Hasse terms in (6) have exponents
$a$ or $a+1$ modulo6, so cannot contribute to $N\equiv b+1$.
Only $P(T)A_0^{a-1}(A_0-a(1+T))$ survives.
Its selected-power degree is at most $a+6<b+7$.
A carry lowers digit sum by six, so only degree $b+1$ contributes.
It is therefore computed exactly by substituting
$T=x+y$, $A_0=-x-Sy$ and extracting $[xy^b]$.
For $b=6$, this same calculation applies for $2\le a\le5$.
At $a=1$ the extra carry contribution is
\[
 -[t^N]C(T)-v[t^N]C(A_0)=-1+vS^6.
\]
Indeed seven selected powers give exactly one copy of1 and six
copies of$L$, with multinomial coefficient7; after division,
$[t^N]C(T)=1$, $[t^N]C(A_0)=a_0a_r^6=-S^6$.
The $a=1$ Artin--Hasse correction is zero here: modulo seven,
$A_0^7$ is supported on powers of seven, and $N=6L+1$ is not
a sum of one or two such powers.

Finally (7) gives $(b+1)[T^{b+1}]P$ for $b\le4$.
For $b=5$ its extra term is $6v[t^N]A_0^6=vS^5$.
For $b=6$ its extra terms are
\[
 [t^N]C(T)+6v[t^N]TA_0^6=1-vS^5(S-1).
\]
The stated table follows by these finite polynomial extractions.
No finite depth $r$ was substituted in the derivation.
\end{proof}

\begin{lemma}[The universal saturated matrix]\label{lem:universal-matrix}
Let $\mathscr D=v\partial_v+(S-1)\partial_S$.
Rows $R_a$ on $1,X,\ldots,X^4$ are determined from (4) by
\[
\begin{split}
 R_a(XQ)&=(6-a+\mathscr D)R_a(Q)+aR_{a-1}(Q),
                                      \quad1\le a<b,\\
 R_0(XQ)&=(6+\mathscr D)R_0(Q)+v\bar V_5(Q).
\end{split}                                                   \tag{8}
\]
For $b=6$ the last term is zero.
Let $\mathcal M_b(S,v)$ be the six-by-five matrix using these
divided rows for $a<b$ and (2) for $a\ge b$, with columns
$Q=1,X,\ldots,X^4$.
Its specialization $S=S_r(v)$ is exactly the reduction of the
actual matrix after the globally seven-divisible rows have each
been divided once by seven.
\end{lemma}
\begin{proof}
Before evaluation, cyclic extraction of
$\mathcal A F=(c-\theta_z)F+zF$ gives
$V_a(XQ)=(c-a-6v\partial_v)V_a(Q)+aV_{a-1}(Q)$ for $a\ge1$.
For row zero the last term is $7vV_5(Q)/5!$.
Divide the rows for $a<b$ by seven and reduce; $c\equiv6$,
$-6\equiv1$ and $5!\equiv1$.
The chain rule is $\mathscr D$, since $vS_r'=S_r-1$.
All operations act on entire polynomial rows before a root is chosen.
In particular no Euler operator is applied to a fixed nonzero fibre.
\end{proof}

\begin{lemma}[A complete finite certificate]\label{lem:finite-certificate}
For each $b=3,4,5,6$, choose polynomials
$q_b,W_b\in\mathbb F_7[S]$, with $q_b$ monic.
At every rank-drop point of $\mathcal M_b(S,v)$ over any extension
field of $\mathbb F_7$, they satisfy
\[
                    q_b(S)=0,\qquad v=W_b(S).            \tag{9}
\]
In $\mathbb F_7[S]/(q_b)$, the recursion
$x_0=1$, $x_{j+1}=1-W_bx_j^7$ satisfies
\[
 x_{P_b}=1,\qquad \gcd(x_j-S,q_b)=1\quad(0\le j<P_b),      \tag{10}
\]
with
\[
\begin{array}{c|rrrr}
 b&3&4&5&6\\\hline
 \deg q_b&15&20&25&30\\
 P_b&90&30&420&36
\end{array}
\]
Hence the actual specialization $\mathcal M_b(S_r(v),v)$ has
rank five at every $v\in k$, for every $r\ge1$.
\end{lemma}
\begin{proof}
The finite polynomial data are retained in
\href{gap_six_scaled_power_plus_one_certificate.json}
{the complete coefficient certificate}.
Its monomials are triples $(\deg_v,\deg_S,\text{coefficient})$.
It specifies all four matrices, their six maximal minors
$\Delta_{b,i}$, the full coefficient vectors of $q_b,W_b$,
and explicit multiplier polynomials with the identities
\[
\begin{aligned}
 q_b(S)&=\sum_{i=0}^5H_{b,i}(S,v)\Delta_{b,i}(S,v),\\
 v-W_b(S)&=\sum_{i=0}^5K_{b,i}(S,v)\Delta_{b,i}(S,v).
\end{aligned}                                                   \tag{11}
\]
There are respectively $67,114,150,219$ nonzero multiplier monomials.
Expansion over $\mathbb F_7[S,v]$ proves both identities, giving (9).
No Groebner basis computation or factorization is an input to the
verifier.

The same file records every state of each orbit in (10).
Polynomial multiplication, remainder and the Euclidean algorithm
verify every transition, every target gcd and the return to1.
This is a finite certificate for the fixed quotient algebras,
not a prefix of the terminal indices.
The return to the initial state makes the sequence periodic,
so every $x_r-S$ is a unit for arbitrary $r\ge0$.

Suppose an actual rank drop existed. By (9), evaluation at its
value of $S$ induces a homomorphism
$\mathbb F_7[S]/(q_b)\to k$, sending $W_b$ to $v$.
Because $S_0=1$, $S_{j+1}=1-vS_j^7$, it sends $x_r$ to
$S_r(v)=S$. It would therefore send the unit $x_r-S$ to zero,
an impossibility. This proves full rank over the whole residue field,
not just at its seven prime-field elements.
\end{proof}

\begin{theorem}[Four terminal families]\label{thm:families}
For every $b\in\{3,4,5,6\}$ and integer $r\ge0$, at terminal
$N=b7^r+1$ the ordinary factorial moments of every nonzero complex
$X(A+BX^6)$ at orders $N-1,N$ cannot both vanish.
\end{theorem}
\begin{proof}
For $r=0$, $N=4,5,6,7$ belong to the Eisenstein terminal family
in \href{PHYSICAL_PRIME_TOTAL_DEGREE_REDUCTION.tex}
{Corollary \texttt{cor:pair} of the physical-prime note}.
Let $r\ge1$ and
$T(t)=\sum_{j=0}^N\binom Nj(N)_{6j}t^j$.
The same note's \texttt{thm:radius} confines any multiple root
to $u=-1/t=7v$ with $v$ a7-adic unit.
The \href{gap_three_hurwitz_structure.tex}{critical-ideal lemma},
\texttt{lem:critical-ideal}, identifies a multiple root with rank
below five for the cyclic auxiliary rows $B_N,\ldots,B_{N+4}$.
For fixed integer shape $N$ these minors are rational polynomials:
apply the complex lemma at an algebraic multiple root and transport
the polynomial vanishing to a7-adic embedding.
No exponential gauge is evaluated on the boundary.

The identity
$B_{N+j}=\Pi_N(\mathcal A)
 \prod_{\ell=N}^{N+j-1}(\mathcal A+6\ell)1$
makes the columns a unit triangular change of the monomials
$Q=1,X,\ldots,X^4$.
The normalizations (1), followed by the global row divisions,
preserve characteristic-zero rank.
If this actual matrix lost rank, reducing a primitive kernel vector
would give a nonzero kernel of
$\mathcal M_b(S_r(\bar v),\bar v)$ over $k$.
Lemma~\ref{lem:finite-certificate} excludes it.
Thus $T$ is squarefree.

For $AB\ne0$ the normalized ordinary moment polynomials of the pair
are $T+(6/N)tT'$ and $T-(1/N)tT'$.
Their determinant in $T,tT'$ is $-7/N\ne0$ and $T(0)=1$;
a common zero would be a multiple root of$T$.
Degenerate nonzero monomials have no zero positive moment.
\end{proof}

\begin{corollary}[Every nonzero leading digit]\label{cor:scaled-family}
For every $1\le b\le6$ and $r\ge0$, ordinary gap-six Strong
Factorial holds at terminal $N=b7^r+1$.
\end{corollary}
\begin{proof}
For $r=0$, all terminals $N=2,\ldots,7$ belong to the native
Eisenstein family, \texttt{cor:pair} of the
\href{PHYSICAL_PRIME_TOTAL_DEGREE_REDUCTION.tex}{physical-prime note}.
Let $r\ge1$. Theorem~\ref{thm:families} supplies $b=3,4,5,6$.
The cases $b=1$ and $b=2$ are, respectively,
\texttt{thm:family} of the
\href{GAP_SIX_THIN_DIGIT_SFC.tex}{thin-digit note} and
\texttt{thm:family} of the
\href{GAP_SIX_TWO_POWER_PLUS_ONE_SFC.tex}{two-power-plus-one note}.
\end{proof}

\begin{corollary}[The refined degree-seven frontier]\label{cor:degree-seven-frontier}
Let
\[
 \mathcal U=\{b7^s:1\le b\le6,\ s\ge0\}
          \cup\{b7^r+1:1\le b\le6,\ r\ge0\},\qquad
 P_j(t)=\sum_{k=0}^j\binom jk(j+1)_{6k}t^k.
\]
Ordinary binomial Strong Factorial through degree seven is equivalent
to $\gcd_{\mathbb Q[t]}(P_{N-1},P_N)=1$ for every terminal
$N\ge6$ outside $\mathcal U$.
Every hypothetical common root on those remaining tests still has
$v_7(-1/t)=1$ at every algebraic embedding.
\end{corollary}
\begin{proof}
The exact degree-seven equivalence and the critical radius are
\texttt{cor:degree-seven-frontier} of the
\href{GAP_SIX_THIN_DIGIT_SFC.tex}{thin-digit note}.
Corollary~\ref{cor:scaled-family} supplies every additional omitted
terminal. Removing already proved tests preserves sufficiency;
necessity follows since the remaining tests are a subset of the
degree-seven assertion.
\end{proof}
The remaining terminals, all-index gap-six and degree-seven
binomial Strong Factorial, and unrestricted FC/SFC remain open.
No every-positive-real-Gamma-shape extension is asserted here.
The written digit and formal-series arguments give every depth $r$.
The exact finite certificate verifies fixed matrix identities and
complete periodic orbits, not an unrestricted assertion from a prefix.
Internal replay and mathematical audit do not imply independent
whole-proof replay, formal verification, external review or priority.
\end{document}
